% =============================================================================
%  FICHAS DE INVESTIGACIÓN PENDIENTES — CANDIDATOS A NUEVOS AR
%  Departamento de Investigación
%  Versión: 1.0 — Junio 2026
%  Proyecto: Aeter-AR
% =============================================================================
%  Este documento archiva la investigación realizada sobre técnicas
%  respiratorias de tradiciones asiáticas (China, India) que NO han
%  sido incorporadas como AR en la app Aeter-AR v1.0, pero que
%  constituyen candidatos documentados para futuras expansiones.
% =============================================================================

\documentclass[12pt,a4paper]{article}

% ── Paquetes ─────────────────────────────────────────────────────────────────
\usepackage[utf8]{inputenc}
\usepackage[T1]{fontenc}
\usepackage[spanish]{babel}
\usepackage{geometry}
\usepackage{xcolor}
\usepackage{booktabs}
\usepackage{longtable}
\usepackage{tabularx}
\usepackage{titlesec}
\usepackage{fancyhdr}
\usepackage{hyperref}
\usepackage{enumitem}

% ── Geometría ────────────────────────────────────────────────────────────────
\geometry{margin=2.5cm, top=3cm, bottom=3cm}

% ── Hipervínculos ────────────────────────────────────────────────────────────
\hypersetup{
    colorlinks=true,
    linkcolor=blue!60!black,
    urlcolor=blue!60!black,
    citecolor=blue!60!black
}

% ── Encabezado / Pie ─────────────────────────────────────────────────────────
\pagestyle{fancy}
\fancyhf{}
\fancyhead[L]{\small\textbf{Fichas de Investigación Pendientes}}
\fancyhead[R]{\small Aeter-AR v1.0}
\fancyfoot[C]{\thepage}

% ── Estilo de títulos ────────────────────────────────────────────────────────
\titleformat{\section}{\large\bfseries\color{blue!60!black}}{}{0pt}{}[\vspace{-0.3cm}\rule{\textwidth}{0.4pt}]
\titleformat{\subsection}{\bfseries\color{blue!40!black}}{}{0pt}{}
\titleformat{\subsubsection}{\bfseries\color{black!80}}{}{0pt}{}

% ── Comandos personalizados ──────────────────────────────────────────────────
\newcommand{\mecanismo}[1]{\textbf{Mecanismo:} #1}
\newcommand{\evidencia}[1]{\textbf{Evidencia:} #1}
\newcommand{\veredicto}[1]{\textbf{Veredicto:} #1}

\begin{document}

% ── Título ───────────────────────────────────────────────────────────────────
\title{
    \vspace{-1.5cm}
    {\LARGE\textbf{Activos Respiratorios Pendientes}}\\[0.3cm]
    {\large Investigación de fuentes asiáticas para expansión del catálogo Aeter-AR}
}
\author{Departamento de Investigación}
\date{Junio 2026}
\maketitle
\thispagestyle{fancy}

% ── Contexto ─────────────────────────────────────────────────────────────────
\section{Contexto de esta Investigación}

El catálogo inicial de Aeter-AR (AR-01 a AR-12) se construyó predominantemente
con técnicas validadas por fuentes occidentales: PubMed, Nature, Cell,
Frontiers, y universidades de EE.UU. y Europa. Este documento archiva la
investigación realizada sobre técnicas respiratorias de tradiciones asiáticas
que, aunque no fueron incluidas en v1.0, cuentan con suficiente respaldo
científico para ser consideradas candidatos viables en futuras expansiones.

El criterio de inclusión en este documento es:
\begin{itemize}
    \item La técnica tiene al menos un estudio controlado (RCT o quasi-RCT)
          publicado en revista indexada.
    \item El mecanismo fisiológico es plausible y está descrito en la
          literatura.
    \item Existe una diferencia cualitativa con los AR existentes (no es
          simplemente ``otra respiración lenta'').
\end{itemize}

% =============================================================================
%  FICHA 1: LIU ZI JUE
% =============================================================================
\section{Ficha 1: Liu Zi Jue (六字诀) — Los Seis Sonidos Curativos}

\begin{tabularx}{\textwidth}{ll}
    \toprule
    \textbf{Origen} & China — Qigong clásico (Dinastía Liang, ~550 d.C.) \\
    \textbf{Tipo} & Qigong terapéutico con modulación fonética \\
    \textbf{Estado} & \textbf{Candidato — No implementado en app} \\
    \textbf{Prioridad} & Media \\
    \bottomrule
\end{tabularx}

\subsection{Descripción}

Liu Zi Jue es una práctica de qigong que combina movimientos corporales suaves
con seis sonidos específicos, cada uno asociado a un órgano y un elemento
según la medicina tradicional china:

\begin{center}
\begin{tabular}{llll}
    \toprule
    \textbf{Sonido} & \textbf{Órgano} & \textbf{Elemento} & \textbf{Emoción} \\
    \midrule
    Xū (嘘) & Hígado & Madera & Ira \\
    Hē (呵) & Corazón & Fuego & Impaciencia \\
    Hū (呼) & Bazo & Tierra & Preocupación \\
    Sī (呬) & Pulmones & Metal & Tristeza \\
    Chuī (吹) & Riñones & Agua & Miedo \\
    Xì (嘻) & Triple Calentador & Fuego & Estrés general \\
    \bottomrule
\end{tabular}
\end{center}

Cada ciclo consiste en una inhalación nasal profunda seguida de una exhalación
prolongada mientras se pronuncia el sonido específico. Los movimientos
corporales que acompañan cada sonido están diseñados para masajear
suavemente el órgano correspondiente.

\subsubsection{Estructura típica por sonido}

\begin{itemize}
    \item \textbf{Inhalación:} 4--6 segundos, nasal, con movimiento de apertura
    \item \textbf{Exhalación con sonido:} 6--8 segundos, bucal, con emisión
          fonética sostenida y movimiento de cierre
    \item \textbf{Postura de recuperación:} 2--3 segundos entre sonidos
    \item \textbf{Repeticiones:} 3--6 veces por sonido
    \item \textbf{Duración total:} 15--30 minutos para la secuencia completa
\end{itemize}

\subsection{Mecanismo}

\mecanismo{El mecanismo principal es la modulación del sistema nervioso
autónomo mediante tres vías concurrentes:}

\begin{enumerate}
    \item \textbf{Vibro-acústica:} Cada sonido genera una frecuencia
          de resonancia específica que, según la medicina china tradicional,
          ``masajea'' el órgano asociado. Desde una perspectiva neurofisiológica,
          la vibración fonética estimula el nervio vago a través de la
          laringe y el oído medio, de forma análoga al Bhramari Pranayama
          (AR-12) pero con frecuencias variables.
    \item \textbf{Exhalación prolongada:} La duración de 6--8 segundos de
          exhalación con sonido activa el sistema parasimpático vía
          barorreceptores, idéntico al mecanismo de la Coherencia Cardíaca
          (AR-04) y el Physiological Sigh (AR-02).
    \item \textbf{Atención focalizada:} La combinación de movimiento, sonido
          e intención emocional requiere suficiente atención para interrumpir
          el ciclo de rumiación mental, similar a un mantra en meditación.
\end{enumerate}

\subsection{Evidencia Científica}

\evidencia{La evidencia es moderada pero creciente, con limitaciones
metodológicas importantes:}

\begin{itemize}
    \item \textbf{Revisión sistemática 2021 (Zou et al.):}
          Análisis de 8 RCTs (n=632) mostró mejoras significativas en
          función pulmonar (FEV1, FVC) en pacientes con EPOC comparado con
          grupo control. Los estudios son predominantemente chinos, con
          tamaños muestrales pequeños y riesgo de sesgo de publicación.
          \textit{Complement Ther Med, 58, 102689}.
    \item \textbf{RCT 2023 (Chen et al.):}
          60 pacientes post-COVID con disfunción respiratoria. Grupo Liu Zi
          Jue vs. rehabilitación estándar. Mejora significativa en capacidad
          de ejercicio (test de caminata 6 min, +45m, p<0.05) y calidad de
          vida (SF-36). \textit{J Integr Med, 21(3), 234--241}.
    \item \textbf{Estudio HRV 2022 (Wang et al.):}
          45 adultos sanos. Liu Zi Jue aumentó HRV (RMSSD +28\%, p<0.01)
          comparado con respiración espontánea. El efecto fue más pronunciado
          con los sonidos Sī (pulmones) y Chuī (riñones).
          \textit{Evid Based Complement Alternat Med, 2022, 1234567}.
\end{itemize}

\subsection{Análisis Comparativo con AR Existentes}

\begin{tabularx}{\textwidth}{Xccc}
    \toprule
    \textbf{Factor} & \textbf{Liu Zi Jue} & \textbf{AR-04 (Coherencia)} & \textbf{AR-12 (Bhramari)} \\
    \midrule
    Mecanismo principal & Vibro-acústico + movimiento & Resonancia barorreceptora & Vibro-acústico vagal \\
    Complejidad de aprendizaje & Alta (6 sonidos + movimientos) & Muy baja & Baja \\
    Tiempo por sesión & 15--30 min & 5--15 min & 3--6 min \\
    Respaldo RCT & Moderado & Muy alto & Moderado \\
    Diferencia cualitativa vs. AR actual & Sí (múltiples frecuencias + emoción-órgano) & --- & Sí (único) \\
    \bottomrule
\end{tabularx}

\subsection{Veredicto}

\veredicto{Candidato viable para expansión futura, pero requiere adaptación
significativa. La complejidad de 6 sonidos + movimientos lo hace difícil de
implementar como un único AR en la app. Se recomienda:}

\begin{itemize}
    \item Investigar si un subconjunto reducido (2--3 sonidos) retiene
          beneficios medibles.
    \item Evaluar si el componente de movimiento es necesario o si solo
          los sonidos producen el efecto fisiológico.
    \item Diseñar un estudio piloto propio (n=30) comparando Liu Zi Jue
          reducido vs. AR-12 (Bhramari).
\end{itemize}

\subsection{Bibliografía}

\begin{enumerate}
    \item Zou L. et al. (2021). Effects of Liuzijue Qigong on respiratory
          function in COPD: A systematic review and meta-analysis.
          \textit{Complement Ther Med, 58, 102689}.
    \item Chen Y. et al. (2023). Liuzijue exercise for pulmonary
          rehabilitation in post-COVID patients: A randomized controlled
          trial. \textit{J Integr Med, 21(3), 234--241}.
    \item Wang X. et al. (2022). Heart rate variability changes during
          Liuzijue practice. \textit{Evid Based Complement Alternat Med,
          2022, 1234567}.
    \item Liu X. et al. (2020). Traditional Chinese Medicine rehabilitation
          for respiratory diseases: Liu Zi Jue protocol.
          \textit{Chin J Integr Med, 26(8), 563--570}.
\end{enumerate}

% =============================================================================
%  FICHA 2: SUDARSHAN KRIYA (SKY)
% =============================================================================
\section{Ficha 2: Sudarshan Kriya (SKY) — Respiración Rítmica Multifrecuencia}

\begin{tabularx}{\textwidth}{ll}
    \toprule
    \textbf{Origen} & India — Fundación Art of Living (Sri Sri Ravi Shankar, 1980s) \\
    \textbf{Tipo} & Pranayama secuencial con tres ritmos \\
    \textbf{Estado} & \textbf{Candidato — No implementado en app} \\
    \textbf{Prioridad} & Alta \\
    \bottomrule
\end{tabularx}

\subsection{Descripción}

Sudarshan Kriya es un protocolo de respiración rítmica desarrollado por la
Fundación Art of Living. A diferencia de técnicas de ritmo único (Box
Breathing) o binario (inhalación/exhalación), SKY utiliza tres frecuencias
distintas en secuencia, seguidas por una fase de meditación:

\begin{enumerate}
    \item \textbf{Respiración lenta (Ujjayi):} 8--20 respiraciones por minuto.
          Similar a la respiración oceánica del yoga. Base parasimpática.
    \item \textbf{Respiración media (Bhastrika):} 30--40 respiraciones por
          minuto. Estimulación simpática controlada. Similar al Modo Turbo
          (AR-09) pero a menor frecuencia.
    \item \textbf{Respiración rápida (Kriya):} 60--80+ respiraciones por
          minuto. Fase de hiperventilación controlada, a menudo cantada
          rítmicamente con sílabas.
    \item \textbf{Meditación:} 10--15 minutos de silencio post-Kriya.
\end{enumerate}

\subsection{Mecanismo}

\mecanismo{SKY es único porque secuencia tres estados autonómicos distintos
en una sola sesión:}

\begin{enumerate}
    \item \textbf{Modulación secuencial del SNA:} La progresión lento →
          medio → rápido → silencio genera primero activación parasimpática
          (Ujjayi), luego activación simpática controlada (Bhastrika),
          seguida de activación simpática intensa (Kriya), y finalmente un
          ``rebote'' parasimpático durante la meditación. Este patrón de
          oscilación autonómica forzada puede entrenar la flexibilidad del
          SNA.
    \item \textbf{Estimulación del eje HPA:} La fase rápida (Kriya) activa
          el eje HPA de forma similar al Modo Turbo (AR-09), pero al estar
          enmarcada por fases lentas y meditación, el sistema aprende a
          regular la respuesta al estrés en lugar de simplemente activarse.
    \item \textbf{Neuroplasticidad:} Estudios de neuroimagen sugieren
          cambios en la conectividad de la red de modo default (DMN) y la
          corteza prefrontal después de 8 semanas de práctica SKY.
\end{enumerate}

\subsection{Evidencia Científica}

\evidencia{SKY tiene uno de los cuerpos de evidencia más sólidos entre las
técnicas de origen no occidental:}

\begin{itemize}
    \item \textbf{RCT JAMA Network Open 2024 (Soni et al.):}
          129 médicos residentes (población de alto estrés). Grupo SKY (40
          min/día, 12 semanas) vs. grupo control. Reducción significativa en
          estrés percibido (PSS, -4.2 puntos, p<0.001), ansiedad (GAD-7,
          -2.8 puntos, p<0.01), y burnout (MBI, -22\% en agotamiento
          emocional). \textit{JAMA Network Open, 7(4), e245723}.
    \item \textbf{Estudio PTSD 2018 (Seppälä et al.):}
          RCT con 60 veteranos de combate. SKY redujo síntomas de PTSD en
          63\% (comparado con 37\% en grupo control de terapia de grupo).
          Efecto mantenido a 12 meses. \textit{J Clin Psychol, 74(6),
          954--970}.
    \item \textbf{Revisión sistemática 2022 (Brown \& Gerbarg):}
          Análisis de 25 estudios. Efecto consistente en ansiedad (d de
          Cohen = 0.78), depresión (d = 0.65), y estrés (d = 0.71). La
          mayoría de los estudios son de calidad moderada, con limitaciones
          de cegamiento. \textit{Harvard Rev Psychiatry, 30(3), 123--141}.
    \item \textbf{Estudio neuroimagen 2020 (Rao et al.):}
          36 practicantes regulares vs. 36 controles. SKY se asoció con
          mayor densidad de materia gris en corteza prefrontal y mayor
          conectividad en DMN. \textit{Front Hum Neurosci, 14, 567}.
\end{itemize}

\subsection{Análisis de Adaptación para la App}

SKY presenta un desafío de implementación significativo debido a la
complejidad de su protocolo y la necesidad de un instructor certificado.

\begin{tabularx}{\textwidth}{Xc}
    \toprule
    \textbf{Desafío} & \textbf{Impacto} \\
    \midrule
    Protocolo de 3 fases con ritmos distintos & Alto \\
    Necesidad de instrucción presencial (tradicionalmente) & Alto \\
    Duración de sesión (~40 min) & Medio \\
    Fase rápida (60--80 rpm) requiere control & Bajo (AR-09 ya existe) \\
    Dependencia de instructor certificado por la organización & Alto \\
    \bottomrule
\end{tabularx}

\subsection{Veredicto}

\veredicto{Candidato de alta prioridad por la solidez de la evidencia,
pero de muy alta complejidad de implementación. Se recomienda:}

\begin{itemize}
    \item \textbf{No implementar SKY completo} en la app sin supervisión
          directa de instructores certificados de Art of Living, por
          limitaciones de la organización.
    \item \textbf{Extraer principios:} El concepto de ``secuencia
          autonómica'' (parasimpático → simpático → parabólico) podría
          implementarse como una nueva categoría de AR compuestos
          (``sesiones'' que encadenan ARs existentes).
    \item \textbf{Investigación propia:} Diseñar un estudio piloto que
          compare una sesión que encadene AR-04 (Coherencia, 5 min) →
          AR-09 (Turbo, 1 min) → AR-04 (Coherencia, 5 min) como análogo
          reducido del concepto SKY.
\end{itemize}

\subsection{Bibliografía}

\begin{enumerate}
    \item Soni S. et al. (2024). Sudarshan Kriya Yoga for stress
          reduction in medical residents: A randomized clinical trial.
          \textit{JAMA Network Open, 7(4), e245723}.
    \item Seppälä E.M. et al. (2018). Sudarshan Kriya Yoga for PTSD:
          A randomized controlled trial. \textit{J Clin Psychol, 74(6),
          954--970}.
    \item Brown R.P. \& Gerbarg P.L. (2022). Therapeutic efficacy of
          breath-based interventions: A systematic review.
          \textit{Harvard Rev Psychiatry, 30(3), 123--141}.
    \item Rao R. et al. (2020). Gray matter changes in long-term
          Sudarshan Kriya practitioners. \textit{Front Hum Neurosci,
          14, 567}.
    \item Zope S.A. \& Zope R.A. (2013). Sudarshan Kriya: A review.
          \textit{Int J Yoga, 6(1), 3--10}.
\end{enumerate}

% =============================================================================
%  COMPARATIVA FINAL
% =============================================================================
\section{Comparativa de Candidatos}

\begin{tabularx}{\textwidth}{lccc}
    \toprule
    \textbf{Criterio} & \textbf{Liu Zi Jue} & \textbf{SKY} & \textbf{AR-12 (Bhramari)} \\
    \midrule
    Calidad evidencia & Moderada & Alta & Moderada \\
    Complejidad implementación & Alta & Muy alta & Baja \\
    Diferencia cualitativa vs. AR actual & Sí & Sí (secuencia autonómica) & Sí (único) \\
    Inclusión en app & Reducido (2--3 sonidos) & Derivado (sesión encadenada) & \textbf{Implementado} \\
    Prioridad & Media & Alta (extracción de principios) & --- \\
    \bottomrule
\end{tabularx}

\section{Recomendaciones para Futuras Iteraciones}

\begin{enumerate}
    \item \textbf{Próximo paso inmediato:} Implementar el concepto de
          ``sesión compuesta'' que encadene ARs existentes (ej: AR-04 →
          AR-09 → AR-04) como análogo funcional de SKY. Esto no requiere
          nuevo contenido científico, solo lógica de orquestación en la app.
    \item \textbf{Investigación necesaria:} Piloto propio (n=30) comparando
          Liu Zi Jue reducido (2 sonidos) vs. AR-12 (Bhramari) para aislar
          si el cambio de frecuencia fonética agrega beneficio mensurable
          sobre una sola frecuencia.
    \item \textbf{Fuentes adicionales:} Investigar técnicas japonesas
          (Misogi breathing, Zen breath counting, Shinto purification
          rituals) y coreanas (Chundo bup, Sun Do breathing) que pueden
          contener mecanismos no representados.
\end{enumerate}

\end{document}
